\documentclass[../Main.tex]{subfiles}

\begin{document}
We will consider always the velocity field $\vec{u}(\vec{x}, t)$ that describes the velocity of individual fluid particles at the point $\vec{x}$ and time $t$.
\section{Streamlines and Pathlines}
\begin{definition}{Streamline}
    A \underline{streamline} is a curve that is everywhere tangential to the flow of the fluid at a given instant in time.
\end{definition}
A streamline through $\vec{x_0}$ at time $t_0$ is given in parametric form as $\vec{x}(s;\vec{x_0}, t_0)$ and solves:
\begin{equation*}
    \frac{d \vec{x}}{d s} = \vec{u}(\vec{x}, t_0),\qquad\vec{x}(0;\vec{x_0}, t_0) = \vec{x_0}
\end{equation*}
\begin{definition}{Pathline}
    A \underline{pathline} is the trajectory of a fluid particle starting at a given point and time.
\end{definition}
The pathline starting at $\vec{x_0}$ at time $t_0$ is given in parametric form $\vec{x}(t;\vec{x_0})$ and solves:
\begin{equation*}
    \frac{d\vec{x}}{dt} = \vec{u}(\vec{x}, t),\qquad\vec{x}(t_0;\vec{x_0}) = \vec{x_0}
\end{equation*}
\begin{definition}{Steady flow}
    A fluid flow is \underline{steady} if its velocity field does not depend on time.
\end{definition}
In the case of steady flow, pathlines and streamlines are the same.
\section{The Material Derivative}
\begin{definition}{Material derivative}
    The \underline{material derivative} of a vector (or scalar) field $\vec{F}$ is given by:
    \begin{equation*}
        \frac{D \vec{F}}{D t} = \frac{\partial \vec{F}}{\partial t} + (\vec{u} \cdot \nabla) \vec{F}
    \end{equation*}
    where the first term is the \underline{local derivative} and the second term is the \underline{advected derivative}.
\end{definition}
\section{Conservation of Mass}
\begin{definition}{Mass density}
    The \underline{mass density} of a fluid is the mass per unit volume $\rho(\vec{x}, t)$ of the fluid.
\end{definition}
Considering the flow of mass through a small volume $\delta V$ and taking the limit as this goes to $0$, we find:
\begin{equation}
    \frac{\partial \rho}{\partial t} + \nabla \cdot (\rho\vec{u}) = 0
    \label{eqnConserveMass}
\end{equation}
\begin{definition}{Incompressible}
    A fluid flow is \underline{incompressible} if $\frac{D \rho}{D t} = 0$. This is equivalent to $\nabla \cdot \vec{u} = 0$.
\end{definition}
We take $\rho$ constant, so $\nabla \cdot \vec{u} = 0$ anyway.
\section{Kinematic Boundary Conditions}
For a surface with normal $\vec{n}$ and velocity $\vec{U}$, the Kinematic Boundary Condition (KBC) is:
\begin{equation}
    \vec{u} \cdot \vec{n} = \vec{U} \cdot \vec{n}.
    \label{eqnKBC}
\end{equation}
For a fixed boundary, we therefore require $\vec{u} \cdot \vec{n} = 0$.

For a free surface defined by $z = \zeta(x, y, t)$ then the KBC becomes:
\begin{equation*}
    \frac{D F}{D t} = 0
\end{equation*}
where $F(x, y, z, t) = z - \zeta(x, y, t)$, and the surface is $F = 0$.
\section{Streamfunctions}
For incompressible flow $\nabla \cdot \vec{u} = 0$ and so we can construct a vector potential $\vec{A}$ such that $\vec{u} = \nabla \times \vec{A}$.
\subsection{2D Cartesian Coordinates}
We consider:
\begin{equation*}
    \vec{u} =
    \begin{pmatrix}u \\ v \\ 0\end{pmatrix}
\end{equation*}
and $\vec{A} = \psi(x, y) \vec{e_z}$. Therefore:
\begin{equation*}
    u = \frac{\partial \psi}{\partial y},\qquad v = -\frac{\partial \psi}{\partial x}
\end{equation*}
\subsection{2D Polar Coordinates}
Now consider a flow $\vec{u} = u_r \vec{e_r} + u_\theta \vec{e_\theta}$. Then $\vec{A} = \phi(r, \theta)\vec{e_z}$ given by:
\begin{equation*}
    \frac1r \frac{\partial \phi}{\partial \theta} = u_r(r,\theta),\qquad -\frac{\partial \phi}{\partial r} = u_\theta(r, \theta)
\end{equation*}
\end{document}